\section{Actualitzacio final de la infraestructura}

La plataforma inicial del projecte estava desplegada en una subscripcio
empresarial externa. Quan aquesta subscripcio va quedar deshabilitada o en mode
nomes lectura, el cluster AKS original va deixar de ser operable. La solucio no
ha estat reparar aquest entorn, sino reconstruir el desplegament en una
subscripcio Azure for Students mantenint el mateix repositori, els mateixos
manifests Kubernetes i el mateix flux d'execucio.

Aquesta migracio es important per al projecte perque demostra que la plataforma
no depen d'un cluster concret. Si es disposa d'una subscripcio activa, un ACR i
un AKS compatible, el sistema es pot tornar a aixecar aplicant els manifests del
repositori i reconstruint les imatges. Les mesures obtingudes despres de la
migracio, pero, s'han d'interpretar com a resultats del nou entorn i no com una
continuacio directa de l'entorn empresarial anterior.

\subsection{Configuracio del cluster}

El cluster final s'ha creat a la regio \texttt{spaincentral}, amb Kubernetes
\texttt{1.34.7}, Azure CNI overlay i node pool \texttt{Standard\_B2s\_v2}. La
fase inicial va funcionar amb un sol node, pero per a la demo final s'ha ampliat
a tres nodes. El motiu es practic: amb un sol node no era estable mantenir
Backstage, Elasticsearch, Kafka, NATS, RabbitMQ i els generadors de carrega al
mateix temps.

El repartiment final ha estat:

\begin{itemize}
  \item \texttt{aks-nodepool1-10848180-vmss000001}: Kafka i Strimzi.
  \item \texttt{aks-nodepool1-10848180-vmss000002}: orquestrador i Jobs
        \texttt{load-generator}, etiquetat amb \texttt{benchmark-role=loadgen}.
  \item \texttt{aks-nodepool1-10848180-vmss000003}: NATS i RabbitMQ.
\end{itemize}

Aquest repartiment no converteix el cluster en un entorn de benchmarking
perfecte, pero redueix el soroll mes greu: evita que els generadors de carrega
caiguin de manera aleatoria sobre nodes diferents i separa, tant com permet la
subscripcio d'estudiant, brokers i carrega.

\subsection{Canvis aplicats durant la migracio}

Durant la reconstruccio s'han detectat i corregit diversos problemes:

\begin{itemize}
  \item La prova Confluent intentava connectar-se a
        \texttt{redpanda.brokers.svc.cluster.local:9093}. En el cluster final
        Confluent es documenta com a cami Kafka-compatible i usa
        \texttt{kafka-cluster-kafka-bootstrap.brokers.svc.cluster.local:9092}.
  \item NATS s'ha estabilitzat amb \texttt{nats://nats.brokers.svc.cluster.local:4222}.
        El servei \texttt{nats-headless} queda disponible per monitoratge i
        comprovacions internes.
  \item RabbitMQ s'ha mantingut via
        \texttt{amqp://admin:<password>@rabbitmq.brokers.svc.cluster.local:5672}.
  \item Els Jobs de benchmark ara es fixen al node
        \texttt{benchmark-role=loadgen}.
  \item L'orquestrador incorpora cua interna i diferencia els estats
        \texttt{pending}, \texttt{running}, \texttt{completed},
        \texttt{failed} i \texttt{cancelled}.
  \item La pantalla de Resultats mostra els runs en curs a la fila principal i
        mou els pendents a una llista separada.
\end{itemize}

\subsection{Concurrencia i qualitat de les dades}

El valor \texttt{MAX\_CONCURRENT\_RUNS} controla quants Jobs poden entrar al
cluster al mateix temps. Per a la demo final s'ha deixat a \texttt{3}, de manera
que es poden crear 16 execucions des del portal sense que Kubernetes rebi 16
pods simultanis al node de carrega. Els runs sobrants queden en estat
\texttt{pending} fins que hi ha espai.

Per a una taula final de benchmarking, el valor recomanat continua sent
\texttt{1}. Amb tres runs simultanis s'obte una demo mes rapida, pero els tres
generadors comparteixen node i poden introduir soroll en latencia i throughput.
Per tant, els resultats de demo serveixen per validar el flux complet; els
resultats que es defensin com a comparacio estricta s'han de repetir en mode
serial o explicitar aquesta limitacio.

\subsection{Cost de l'entorn}

Azure for Students proporciona 100 USD de credit durant 12 mesos i no requereix
targeta de credit. Aquest credit cobreix consum real, pero no elimina el cost
dels recursos. AKS pot estar en nivell de gestio Free, pero els nodes, discs i
registre d'imatges consumeixen credit.

Segons Azure Retail Prices API consultada el 23/05/2026 per a
\texttt{spaincentral}, una VM \texttt{Standard\_B2s\_v2} Linux costa
aproximadament 0.0912 USD/h. Amb tres nodes:

\[
3 \times 0.0912 = 0.2736 \text{ USD/h}
\]

Si el cluster estigues encès 720 hores al mes:

\[
0.2736 \times 720 = 196.99 \text{ USD/mes}
\]

A aquesta xifra cal afegir aproximadament 5 USD/mes d'Azure Container Registry
Basic, l'emmagatzematge d'imatges, els discs gestionats i possibles costos de
transit. En consequencia, el cluster no s'ha de mantenir encès permanentment:
amb tres nodes pot consumir el credit estudiantil en menys d'un mes si no
s'apaga.

Sense Azure for Students, el cost tecnic seria similar en tarifa Pay-As-You-Go,
pero es facturaria directament. Amb Azure for Students, el mateix consum resta
del credit academic disponible.

\subsection{Estat final de l'aplicacio}

L'aplicacio final queda organitzada en cinc pantalles principals. Home introdueix
els conceptes de productor, broker i consumidor. Cataleg documenta
arquitectures, protocols i plataformes amb versions i notes de reproduibilitat.
Escenaris permet crear i executar combinacions. Execucions mostra l'historial i
els estats. Resultats mostra les metriques en directe i l'historic.

El cataleg final inclou 15 components predefinits, inclosa l'arquitectura
\texttt{SEA}. Les fitxes documenten la versio coneguda dels brokers i deixen
clar que Confluent s'avalua pel cami Kafka-compatible del cluster, no com una
instal.lacio completa de Confluent Platform amb Schema Registry, ksqlDB o
Control Center.

Les fonts consultades per aquesta seccio son:
\begin{itemize}
  \item Azure for Students: \url{https://azure.microsoft.com/en-us/free/students}
  \item AKS Free, Standard and Premium tiers:
        \url{https://learn.microsoft.com/azure/aks/free-standard-pricing-tiers}
  \item Azure Retail Prices API:
        \url{https://learn.microsoft.com/rest/api/cost-management/retail-prices/azure-retail-prices}
\end{itemize}
